\chapter{Model: linear bandits, identification algorithms and the Gaussian width}
\label{chap:model}

This chapter fixes the objects of the paper (Section~1 and the notations of Section~1.1),
in the form in which they will be formalized on top of the LML framework of sequential
algorithms and environments (\texttt{Learning.Algorithm}, \texttt{Learning.Environment},
\texttt{Learning.IsAlgEnvSeq}, \texttt{Learning.stationaryEnv}).

\section{Action sets and the linear Gaussian environment}
\label{sec:model_env}

\begin{definition}[Action set]\label{def:arm_set}
Let $d \ge 1$. An \emph{action set} (or arm set) is a nonempty compact subset $\Xs \subset \R^d$.
It is \emph{spanning} if $\Span(\Xs) = \R^d$. The paper assumes spanning throughout; in this
blueprint the spanning assumption is stated explicitly in every statement that needs it (the
existence of positive definite designs, the Gaussian width, the upper bounds), and the
definitions of environments, algorithms, regret and PAC guarantees below make sense for every
action set, spanning or not (the structured sets of Chapter~\ref{chap:log_gains} are not
spanning). We write $\ball_d := \{x \in \R^d : \norm{x} \le 1\}$ for the unit
Euclidean ball and $\sphere^{d-1}$ for the unit sphere. For $\theta \in \R^d$ the
\emph{value} of an arm $x$ is $\ip{x}{\theta}$ and an \emph{optimal arm} is
$x_\star(\theta) \in \argmax_{x \in \Xs} \ip{x}{\theta}$ (it exists by compactness).
\end{definition}

\textbf{Lean remark.} The ambient space is \texttt{EuclideanSpace ℝ (Fin d)}. The action
type of the LML algorithms is the subtype $\Xs$ with its Borel $\sigma$-algebra, so that
algorithms play in $\Xs$ by typing. Compactness is only used through the existence of
maximizers of continuous functions and the finiteness of $\sup_{x \in \Xs}\norm{x}$; the
spanning assumption is used through the existence of positive definite design matrices
(Lemma~\ref{lem:exists_posdef_design}) and is a separate hypothesis
(\texttt{Submodule.span ℝ 𝒳 = ⊤}) of the statements that need it. Finite action sets are the important special case,
and several statements are first proved for finite $\Xs$ and then extended by density
(Lemma~\ref{lem:sup_countable_dense}).

\begin{definition}[Linear Gaussian environment]\label{def:env}
\uses{def:arm_set}
For $\theta \in \R^d$, the \emph{linear Gaussian environment} $\nu_\theta$ is the stationary
environment (LML \texttt{stationaryEnv}) with reward kernel
$x \mapsto \Normal(\ip{x}{\theta}, 1)$, $x \in \Xs$: at each round $t$ the learner plays
$x_t \in \Xs$ and observes $y_t = \ip{x_t}{\theta} + \eta_t$ where the $\eta_t$ are i.i.d.\
standard Gaussian, independent of the past and of the learner's randomness.
\end{definition}

\begin{lemma}[Noise representation]\label{lem:noise_representation}
\uses{def:env}
Let $(x_t, y_t)_{t < T}$ be an algorithm-environment sequence for an algorithm with actions
in $\Xs$ and the environment $\nu_\theta$, on $(\Omega, \Prob)$. Then
$\eta_t := y_t - \ip{x_t}{\theta}$, $t < T$, are i.i.d.\ $\Normal(0,1)$ and, for every $t$,
$\eta_t$ is independent of the pair $\big((x_s, y_s)_{s < t},\, x_t\big)$ (jointly, i.e.\ of the
$\sigma$-algebra generated by the history up to $t-1$ together with $x_t$).
\end{lemma}

\begin{proof}
\uses{def:env}
The LML statement \texttt{IsAlgEnvSeq.hasCondDistrib\_feedback} gives that the conditional law
of $y_t$ given the history up to $t-1$ and $x_t$ is $\Normal(\ip{x_t}{\theta},1)$; shifting by
the measurable function $\ip{x_t}{\theta}$ of the conditioning variables gives that the
conditional law of $\eta_t$ is $\Normal(0,1)$, a constant kernel, which is the claimed
independence. Joint independence of the $\eta_t$ follows by induction on $t$.
\end{proof}

\section{Identification algorithms and PAC guarantees}
\label{sec:model_pac}

\begin{definition}[Identification algorithm with budget $T$]\label{def:bai_alg}
\uses{def:env}
An \emph{identification algorithm with budget $T \ge 0$} on $\Xs$ is a pair
$\mathsf{Alg} = (\mathrm{alg}, \rho)$ of an LML algorithm $\mathrm{alg}$ with actions in $\Xs$
and observations in $\R$ (a sampling rule), and a \emph{recommendation rule}: a Markov kernel
$\rho$ from histories of length $T$ (the sequence $(x_t, y_t)_{t<T}$) to $\Xs$. Run against
$\nu_\theta$, it produces the history $\hist_T = (x_t,y_t)_{t<T}$ and then the recommendation
$\rec \sim \rho(\hist_T)$. We write $\Prob_\theta$ for the law of $(\hist_T, \rec)$ (unique by
\texttt{Learning.isAlgEnvSeq\_unique}).
\end{definition}

\textbf{Lean remark.} The randomness of the recommendation is allowed on purpose (the lower
bounds are stated for randomized procedures). Budget $T = 0$ is allowed: the history of length
$0$ lives in a one-point space (\texttt{Fin 0 → 𝒳 × ℝ}) and the recommendation kernel is then
just a probability measure on $\Xs$; this degenerate case is needed by the singular-design lemma
(Lemma~\ref{lem:singular_design_fails}) and its use in Chapter~\ref{chap:separation}. Histories
of length $T \ge 1$ are LML's \texttt{Iic (T-1) → 𝒳 × ℝ}. A deterministic recommendation is the kernel
\texttt{Kernel.deterministic}. Stopping times are not modelled: the paper's PAC algorithms use
a deterministic budget $T = (H_1 \log(1/\delta) + H_2)/\eps^2$ and all its lower bounds are
stated for such budgets; a randomized budget can be handled by padding
(Remark~\ref{rem:stopping_times}).

\begin{definition}[Simple regret]\label{def:regret}
\uses{def:arm_set}
The \emph{simple regret} of $\rec \in \Xs$ for the reward vector $\theta$ is
\[
  r(\rec, \theta) := \max_{x \in \Xs} \ip{x}{\theta} - \ip{\rec}{\theta} \;\ge 0 .
\]
We also set $Z(\theta) := \max_{x, y \in \Xs} \ip{x - y}{\theta}$, so that
$r(\rec,\theta) \le Z(\theta)$ for every $\rec \in \Xs$.
\end{definition}

\begin{definition}[$(\eps,\delta)$-PAC algorithm]\label{def:pac}
\uses{def:bai_alg, def:regret}
Let $\eps > 0$ and $\delta \in (0,1)$. An identification algorithm with budget $T$ is
\emph{$(\eps,\delta)$-PAC on $\Xs$} if for every $\theta \in \R^d$,
\[
  \Prob_\theta\big( r(\rec, \theta) \le \eps \big) \ge 1 - \delta .
\]
\end{definition}

\begin{lemma}[Expected regret of a PAC algorithm under a prior]\label{lem:pac_expected_regret}
\uses{def:pac}
Let $\mathsf{Alg}$ be $(\eps,\delta)$-PAC with budget $T$ and let $\pi$ be a probability
measure on $\R^d$ such that $\int Z(\theta)\,\pi(d\theta) < \infty$. Let $\kappa$ be a Markov
kernel from $\R^d$ to the space of pairs $(\hist_T, \rec)$ such that $\kappa(\theta) = \Prob_\theta$
for $\pi$-almost every $\theta$ (``run $\mathsf{Alg}$ against $\nu_\theta$ conditionally on
$\theta$''), and let $(\theta, \hist_T, \rec) \sim \pi \otimes \kappa$. Then
\[
  \E\big[ r(\rec, \theta) \big] \le \eps + \delta \int Z(\theta)\, \pi(d\theta) .
\]
\end{lemma}

\begin{proof}
\uses{def:regret}
Pathwise, $r \le \eps + Z(\theta)\indic\{r > \eps\}$. Taking expectations under
$\pi \otimes \kappa$ (Fubini for the composition-product, Mathlib
\texttt{MeasureTheory.Measure.integral\_compProd}),
$\E[Z(\theta)\indic\{r>\eps\}] = \int Z(\theta)\,\kappa(\theta)(r > \eps)\,\pi(d\theta)
= \int Z(\theta)\,\Prob_\theta(r > \eps)\,\pi(d\theta) \le \delta \int Z\,d\pi$.

\textbf{Lean remark.} The kernel $\kappa$ is a hypothesis, not a construction: measurability of
$\theta \mapsto \Prob_\theta$ for LML trajectory measures is not available off the shelf. Both
uses of the lemma supply $\kappa$ explicitly: for a fixed design
$\kappa(\theta) = \Normal(X\theta, I_T) \otimes \rho$ (Lemma~\ref{lem:fixed_design_law}), and
Chapter~\ref{chap:pre_bayes} builds $\kappa$ from the likelihood ratio of
Lemma~\ref{lem:env_likelihood_ratio}.
\end{proof}

\begin{remark}[Stopping times]\label{rem:stopping_times}
An algorithm with a stopping time $\tau$ and $\E[\tau] \le T_0$ can be turned into one with
deterministic budget $T = T_0/\delta$ by padding with dummy rounds: by Markov's inequality
$\Prob(\tau > T) \le \delta$, so an $(\eps,\delta)$-PAC algorithm with random budget becomes
$(\eps, 2\delta)$-PAC with budget $T_0/\delta$. All lower bounds below therefore transfer to
expected budgets up to the factor $1/\delta$, which is what the paper does in
Appendix~I. This reduction is not part of the formalization plan.
\end{remark}

\section{Non-adaptive fixed designs}
\label{sec:model_fixed}

\begin{definition}[Fixed-design algorithm]\label{def:fixed_design}
\uses{def:bai_alg}
A \emph{fixed design} of length $T$ is a deterministic sequence $x_0, \dots, x_{T-1} \in \Xs$.
The associated \emph{fixed-design (non-adaptive) algorithm} is the identification algorithm
whose sampling rule is deterministic and ignores the observations, playing $x_t$ at round
$t$ (LML \texttt{detAlgorithm} with a history-independent next-action function), together with
an arbitrary recommendation kernel $\rho$ from $(y_t)_{t<T} \in \R^T$ to $\Xs$ (budget
$T \ge 0$; for $T = 0$ the design is empty and $\rho$ is a probability measure on $\Xs$). We write
$X \in \R^{T \times d}$ for the matrix with rows $x_t^\top$ and
$\Sigma_T := X^\top X = \sum_{t<T} x_t x_t^\top$ for its \emph{design matrix}. Since the actions
are deterministic, the history $\hist_T$ is the image of $y = (y_t)_{t<T}$ under the measurable
bijection $y \mapsto ((x_t, y_t))_{t<T}$; we identify $(\hist_T, \rec)$ with $(y, \rec)$ and
write $\Prob_\theta$ for the law of either.
\end{definition}

\begin{lemma}[Law of the observations under a fixed design]\label{lem:fixed_design_law}
\uses{def:fixed_design, lem:noise_representation}
Under a fixed design run against $\nu_\theta$, the observation vector
$y = (y_t)_{t<T}$ satisfies $y = X\theta + \eta$ with $\eta \sim \Normal(0, I_T)$, and
the recommendation is $\rec \sim \rho(y)$. In particular the law of $(y, \rec)$ is
$\Normal(X\theta, I_T) \otimes \rho$.
\end{lemma}

\begin{proof}
\uses{lem:noise_representation, lem:gauss_pi}
Lemma~\ref{lem:noise_representation} with deterministic actions: the $\eta_t$ are i.i.d.\
standard Gaussians, so $\eta \sim \Normal(0,I_T)$ (Lemma~\ref{lem:gauss_pi}).
\end{proof}

\section{Design matrices and the Gaussian width}
\label{sec:model_width}

\begin{definition}[Design distributions and design matrices]\label{def:design_matrix}
\uses{def:arm_set}
A \emph{design distribution} on $\Xs$ is a finitely supported probability vector
$\lambda = (\lambda_x)_{x \in S}$, $S \subset \Xs$ finite, $\lambda_x \ge 0$,
$\sum_{x \in S} \lambda_x = 1$; we write $\lambda \in \simplex_{\Xs}$ and
$\supp(\lambda) := \{x \in S : \lambda_x > 0\}$. Its \emph{design matrix} is
\[
  A(\lambda) := \sum_{x \in S} \lambda_x\, x x^\top \in \PSD .
\]
The \emph{set of design matrices} is $\designset(\Xs) := \{A(\lambda) : \lambda \in
\simplex_{\Xs}\} = \conv\{x x^\top : x \in \Xs\}$. Finally, for a fixed design
$x_0,\dots,x_{T-1}$ the \emph{empirical design distribution} is
$\lambda_T := \frac1T \sum_{t<T} \delta_{x_t}$, so that $\Sigma_T = T\, A(\lambda_T)$.
\end{definition}

\textbf{Lean remark.} A design distribution is a finitely supported function
$\Xs \to [0,1]$ summing to $1$ (\texttt{Finsupp}, or a \texttt{Finset} with a weight function;
zero weights are allowed and $\supp$ is the set of positive weights). Restricting to finitely
supported designs loses nothing: by
Carathéodory's theorem (Lemma~\ref{lem:caratheodory_design}) every barycenter of
$\{xx^\top : x \in \Xs\}$ under a probability measure on $\Xs$ is also $A(\lambda)$ for a
$\lambda$ supported on at most $d(d+1)/2 + 1$ points. Finitely supported designs are what
the fixed-design construction consumes.

\begin{lemma}[Basic properties of the set of design matrices]\label{lem:designset_basic}
\uses{def:design_matrix}
$\designset(\Xs)$ is a convex compact subset of the symmetric matrices, contained in $\PSD$,
and every $A \in \designset(\Xs)$ satisfies $\tr(A) \le \sup_{x\in\Xs}\norm{x}^2$.
\end{lemma}

\begin{proof}
\uses{def:design_matrix}
The map $x \mapsto xx^\top$ is continuous, so $\{xx^\top : x\in\Xs\}$ is compact; the convex
hull of a compact subset of a finite-dimensional space is compact
(Mathlib \texttt{Set.Finite.isCompact\_convexHull} / \texttt{isCompact\_convexHull}).
Positive semidefiniteness and the trace bound are preserved by convex combinations.
\end{proof}

\begin{lemma}[Existence of a positive definite design]\label{lem:exists_posdef_design}
\uses{def:design_matrix}
If $\Xs$ is spanning, there exists $\lambda \in \simplex_{\Xs}$ with $A(\lambda)$ positive
definite. More precisely, if $b_1, \dots, b_d \in \Xs$ is a basis of $\R^d$ (which exists since
$\Span(\Xs) = \R^d$),
then the uniform design on $\{b_1,\dots,b_d\}$ has $A(\lambda) \succ 0$.
\end{lemma}

\begin{proof}
\uses{def:design_matrix}
For $v \ne 0$, $v^\top A(\lambda) v = \frac1d \sum_i \ip{b_i}{v}^2 > 0$ since some
$\ip{b_i}{v} \neq 0$ (the $b_i$ span $\R^d$).
\end{proof}

\begin{definition}[Gaussian width of a compact set for a matrix]\label{def:width_matrix}
\uses{def:arm_set}
For a nonempty compact set $K \subseteq \R^d$ (typically $K = \Xs$, or a region $K = R \subseteq \Xs$
in Chapter~\ref{chap:log_gains}), a positive definite matrix $A \in \PD$ and
$g \sim \Normal(0, I_d)$, the \emph{Gaussian width of $K$ for $A$} is
\[
  \gwidth{K}{A} := \E\Big[ \sup_{x \in K} \ip{x}{A^{-1/2} g} \Big]
  = \E\Big[ \sup_{x \in K} \ip{A^{-1/2}x}{g} \Big].
\]
Here $A^{-1/2}$ is the positive definite square root of $A^{-1}$ (Mathlib \texttt{CFC.sqrt}).
The lemmas of this section that do not involve the design set $\designset(\Xs)$
(Lemmas~\ref{lem:width_matrix_wd}, \ref{lem:width_homog}, \ref{lem:width_mono} and
\ref{lem:width_symm}) are stated and proved for an arbitrary nonempty compact index set $K$;
they are written with $K = \Xs$ for readability.
\end{definition}

\begin{lemma}[The width for a matrix is well defined]\label{lem:width_matrix_wd}
\uses{def:width_matrix, lem:sup_countable_dense}
(For an arbitrary nonempty compact index set $\Xs$.) The random variable $\sup_{x\in\Xs}\ip{A^{-1/2}x}{g}$ is measurable, integrable, and
$0 \le \gwidth{\Xs}{A} \le \sup_{x\in\Xs}\norm{A^{-1/2}x}\cdot \sqrt d$.
Moreover $\gwidth{\Xs}{A}$ only depends on the law of the Gaussian vector
$G := A^{-1/2}g \sim \Normal(0, A^{-1})$: $\gwidth{\Xs}{A} = \E[\sup_{x\in\Xs}\ip{x}{G}]$ for
any $G \sim \Normal(0,A^{-1})$.
\end{lemma}

\begin{proof}
\uses{lem:sup_countable_dense, lem:gauss_norm_moments, lem:gauss_linear_image}
Measurability: the supremum can be taken over a countable dense subset of $\Xs$
(Lemma~\ref{lem:sup_countable_dense}). Integrability and the upper bound:
$\abs{\sup_x \ip{A^{-1/2}x}{g}} \le \sup_x\norm{A^{-1/2}x}\,\norm{g}$ and
$\E\norm{g} \le \sqrt{\E\norm{g}^2} = \sqrt d$ (Lemma~\ref{lem:gauss_norm_moments}).
Nonnegativity: for any $x_0\in\Xs$, $\sup_x \ip{A^{-1/2}x}{g} \ge \ip{A^{-1/2}x_0}{g}$, whose
expectation is $0$. The last claim is that $A^{-1/2}g \sim \Normal(0,A^{-1})$
(Lemma~\ref{lem:gauss_linear_image}) and that the expectation of a measurable function only
depends on the law.
\end{proof}

\begin{definition}[Gaussian width of $\Xs$]\label{def:width}
\uses{def:width_matrix, def:design_matrix, lem:exists_posdef_design}
Let $\Xs$ be spanning. The \emph{Gaussian width term} of the action set is
\[
  \gw(\Xs) := \inf\big\{ \gwidth{\Xs}{A} : A \in \designset(\Xs) \cap \PD \big\}
  = \inf_{\lambda \in \simplex_{\Xs},\ A(\lambda)\succ 0}
    \E\Big[\max_{x\in\Xs}\ip{x}{A(\lambda)^{-1/2}g}\Big].
\]
The set over which the infimum is taken is nonempty by Lemma~\ref{lem:exists_posdef_design},
so $0 \le \gw(\Xs) < \infty$.
\end{definition}

\begin{lemma}[Homogeneity]\label{lem:width_homog}
\uses{def:width_matrix, lem:sqrt_props}
(For an arbitrary nonempty compact index set $\Xs$.) For $A \in \PD$ and $c > 0$,
$\gwidth{\Xs}{cA} = c^{-1/2}\,\gwidth{\Xs}{A}$.
\end{lemma}

\begin{proof}
\uses{def:width_matrix}
$(cA)^{-1/2} = c^{-1/2}A^{-1/2}$ and the supremum is positively homogeneous.
\end{proof}

\begin{lemma}[Monotonicity in the Loewner order]\label{lem:width_mono}
\uses{def:width_matrix, lem:gauss_comparison, lem:loewner_inv}
(For an arbitrary nonempty compact index set $\Xs$.) If $A, B \in \PD$ and $A \preceq B$, then $\gwidth{\Xs}{B} \le \gwidth{\Xs}{A}$.
\end{lemma}

\begin{proof}
\uses{lem:loewner_inv, lem:gauss_comparison, lem:width_matrix_wd}
By Lemma~\ref{lem:loewner_inv}, $B^{-1} \preceq A^{-1}$. The processes
$x \mapsto \ip{x}{G_B}$ and $x\mapsto\ip{x}{G_A}$ with $G_B\sim\Normal(0,B^{-1})$,
$G_A \sim \Normal(0,A^{-1})$ are centered linear Gaussian processes whose covariance matrices
are ordered, so Lemma~\ref{lem:gauss_comparison} (covariance domination) gives
$\E\sup_x\ip{x}{G_B} \le \E\sup_x\ip{x}{G_A}$; conclude with Lemma~\ref{lem:width_matrix_wd}.
\end{proof}

\begin{lemma}[Invariance under linear isomorphisms]\label{lem:width_iso}
\uses{def:width, lem:width_matrix_wd, lem:gauss_linear_image, lem:loewner_congruence}
Let $\Xs$ be spanning and let $M \in \R^{d\times d}$ be invertible and $\Xs' := M\Xs$. Then for every $A \in \PD$,
$\gwidth{\Xs'}{MAM^\top} = \gwidth{\Xs}{A}$, the map $A \mapsto MAM^\top$ is a bijection
from $\designset(\Xs)\cap\PD$ onto $\designset(\Xs')\cap\PD$, and $\gw(\Xs') = \gw(\Xs)$.
\end{lemma}

\begin{proof}
\uses{lem:gauss_law_of_cov, lem:width_matrix_wd}
The process $x \mapsto \ip{Mx}{G'}$ with $G' \sim \Normal(0,(MAM^\top)^{-1})$ is centered
Gaussian with covariance $x^\top M^\top (MAM^\top)^{-1} M y = x^\top A^{-1} y$, the same as
$x \mapsto \ip{x}{G}$, $G \sim \Normal(0,A^{-1})$. Two centered Gaussian vectors with the same
covariance have the same law (Lemma~\ref{lem:gauss_law_of_cov}), so the two suprema have the
same expectation by Lemma~\ref{lem:width_matrix_wd}. The design matrix of the pushed-forward
design is $A(M\lambda) = \sum \lambda_x (Mx)(Mx)^\top = MA(\lambda)M^\top$, and congruence by
an invertible matrix preserves positive definiteness.
\end{proof}

\begin{lemma}[Width of an empirical design]\label{lem:width_empirical}
\uses{def:width, def:design_matrix, lem:width_homog}
Let $\Xs$ be spanning and let $x_0,\dots,x_{T-1} \in \Xs$ with $\Sigma_T = \sum_{t<T} x_tx_t^\top$ positive definite.
Then $\gwidth{\Xs}{\Sigma_T} \ge \gw(\Xs)/\sqrt T$.
\end{lemma}

\begin{proof}
\uses{lem:width_homog, def:width}
$\Sigma_T = T A(\lambda_T)$ with $A(\lambda_T) \in \designset(\Xs)\cap\PD$, so
$\gwidth{\Xs}{\Sigma_T} = T^{-1/2}\gwidth{\Xs}{A(\lambda_T)} \ge T^{-1/2}\gw(\Xs)$.
\end{proof}

\begin{lemma}[Symmetry of the width]\label{lem:width_symm}
\uses{def:width_matrix}
(For an arbitrary nonempty compact index set $\Xs$.) For $A \in \PD$ and $G \sim \Normal(0,A^{-1})$,
$\E[\max_{x\in\Xs}\ip{x}{-G}] = \E[\max_{x\in\Xs}\ip{x}{G}] = \gwidth{\Xs}{A}$ and
$\E[Z(G)] = \E[\max_{x,y\in\Xs}\ip{x-y}{G}] = 2\,\gwidth{\Xs}{A}$.
\end{lemma}

\begin{proof}
\uses{lem:gauss_law_of_cov}
$-G \sim \Normal(0,A^{-1})$ as well (Lemma~\ref{lem:gauss_law_of_cov}), and
$\max_{x,y}\ip{x-y}{G} = \max_x\ip{x}{G} + \max_y\ip{y}{-G}$.
\end{proof}

\section{Sequential composition of algorithms}
\label{sec:model_composition}

The adaptive algorithms of the paper (Chapters~\ref{chap:log_gains},~\ref{chap:norm}
and~\ref{chap:separation}) are built by running a first algorithm for $T_1$ rounds and then
a second algorithm whose parameters depend on the first history. The following framework
lemma isolates what is needed.

\begin{lemma}[Sequential composition]\label{lem:composition}
\uses{def:bai_alg}
Let $\mathrm{alg}_1$ be an algorithm with actions in $\Xs$, $T_1 \ge 1$, and for each history
$h$ of length $T_1$ let $\mathrm{alg}_2(h)$ be an algorithm, measurably in $h$. There is an
algorithm $\mathrm{alg}$ (the \emph{composition}) whose policy coincides with that of
$\mathrm{alg}_1$ at rounds $t < T_1$ and with that of $\mathrm{alg}_2(\hist_{T_1})$ at rounds
$t \ge T_1$ (applied to the history after time $T_1$). For every environment $\nu$, under the
law of the composition against $\nu$: the history $\hist_{T_1}$ has the law of
$\mathrm{alg}_1$ against $\nu$, and conditionally on $\hist_{T_1} = h$ the shifted process
$(x_{T_1+s}, y_{T_1+s})_{s\ge 0}$ is an algorithm-environment sequence for
$(\mathrm{alg}_2(h), \nu)$. Consequently, if $(E_h)_h$ is a family of sets of trajectories, jointly measurable in
$(h, \cdot)$, with
$\Prob^{\mathrm{alg}_2(h),\nu}(E_h) \ge 1-\delta$ for every $h$ in a set of $\mathrm{alg}_1$-probability
at least $1 - \delta'$, then the composition satisfies
$\Prob\big((x_{T_1+s},y_{T_1+s})_{s \ge 0} \in E_{\hist_{T_1}}\big) \ge 1 - \delta - \delta'$.
\end{lemma}

\begin{proof}
\uses{def:bai_alg}
Define the policy of $\mathrm{alg}$ at time $t\ge T_1$ from the history $h'$ of length $t$ by
splitting $h' = (h, h'')$ with $h$ of length $T_1$ and applying the policy of
$\mathrm{alg}_2(h)$ at time $t - T_1$ to $h''$ (with the initial action distribution of
$\mathrm{alg}_2(h)$ at $t = T_1$); measurability in $h'$ holds by assumption. The trajectory
measure of the composition (\texttt{Learning.trajMeasure}, an Ionescu--Tulcea construction)
is, by construction of the kernels, the composition-product of the law of $\hist_{T_1}$ under
$\mathrm{alg}_1$ with the kernel $h \mapsto$ trajectory measure of $(\mathrm{alg}_2(h),\nu)$
shifted by $T_1$. The probability bound is then Fubini for the composition-product.
\end{proof}

\textbf{Lean remark.} This lemma is a general statement about the LML framework, to be
contributed to LML. The measurability of $h \mapsto \mathrm{alg}_2(h)$ means that the policy
kernels and initial distributions, as functions of $(h, \cdot)$, are kernels; in all uses
below $\mathrm{alg}_2(h)$ is a fixed algorithm reparametrized by a measurable function of $h$
(a finite set of candidate arms, or a scale $r_0$).
